% 图契约：本题的源、边界和推导逐项对应；箭头符号按图注解释。
\documentclass[tikz,border=7pt]{standalone}
\usepackage{fontspec,xeCJK,amsmath,amssymb}
\IfFontExistsTF{Times New Roman}{\setmainfont{Times New Roman}}{\setmainfont{TeX Gyre Termes}}
\IfFontExistsTF{Heiti SC}{\setCJKmainfont[AutoFakeBold=1.5]{Heiti SC}}{\setCJKmainfont[AutoFakeBold=1.5]{FandolHei-Regular}}
\usetikzlibrary{arrows.meta,bending,calc,patterns,positioning}
\definecolor{ink}{HTML}{243745}
\definecolor{blue}{HTML}{275DAB}
\definecolor{green}{HTML}{177C63}
\definecolor{red}{HTML}{BD4935}
\definecolor{orange}{HTML}{AA6C16}
\definecolor{purple}{HTML}{715393}
\tikzset{every node/.style={font=\fontsize{9.3}{12.5}\selectfont,text=ink},
axis/.style={-{Stealth[length=4pt]},line width=.65pt,draw=ink!65},
vec/.style={-{Stealth[length=5pt]},line width=1pt,draw=blue},
short/.style={-{Stealth[length=3pt]},line width=.8pt},
guide/.style={dashed,draw=ink!40,line width=.5pt},
card/.style={rounded corners=2pt,draw=blue!65,fill=blue!4,align=center,text width=6.5cm,inner xsep=9pt,inner ysep=7pt},
title/.style={font=\fontsize{10.5}{14}\selectfont\bfseries,text width=15cm,align=center},
label/.style={font=\fontsize{8.4}{11}\selectfont,align=center},
note/.style={label,text width=6.2cm}}
\begin{document}
\begin{tikzpicture}[x=1cm,y=1cm]
\path[use as bounding box] (-8,-3.5) rectangle(8,3.5);
\node[title] at (0,3.15) {M3-002\quad 由平均分子偶极矩求极化强度};
\begin{scope}[shift={(-4,0)}]
\filldraw[draw=ink!55,fill=blue!4] (-1.8,-1.3) rectangle(1.2,1.3);
\draw[ink!55] (-1.8,1.3)--(-1.2,1.85)--(1.8,1.85)--(1.2,1.3) (1.8,1.85)--(1.8,-.75)--(1.2,-1.3);
\foreach \x in {-1.3,-.4,.5}{\foreach \y in {-.85,0,.85}{\draw[vec,green] (\x,\y)--++(.55,0);}}
\node[label] at(0,-1.8){$N\Delta V$ 个分子；箭头为平均偶极矩};
\node at(0,2.25){宏观小体积 $\Delta V$};
\draw[axis] (-1.7,-2.3)--(1.7,-2.3) node[right] {$x$};\end{scope}
\node[card,draw=blue!65,fill=blue!4] at (3.9,1.1) {$\displaystyle\mathbf P=\frac{\sum\mathbf p_i}{\Delta V}=N\langle\mathbf p\rangle$};
\node[card,draw=green!65,fill=green!4] at (3.9,-1.1) {$\mathbf P=6.0\times10^{-2}\mathbf e_x\ \mathrm{C/m^2}$};
\end{tikzpicture}\end{document}